% Generated draft: requires mathematical and source review.
\documentclass[11pt]{article}
\usepackage[margin=25mm]{geometry}
\usepackage{fontspec}
\setmainfont{Latin Modern Roman}
\usepackage{amsmath,amssymb,mathtools,unicode-math}
\setmathfont{Latin Modern Math}
\usepackage{enumitem,needspace,xurl,hyperref}
\setlength{\parindent}{0pt}
\setlength{\parskip}{5pt}
\setlength{\emergencystretch}{3em}
\newcommand{\N}{\mathbb{N}}
\newcommand{\Z}{\mathbb{Z}}
\newcommand{\Q}{\mathbb{Q}}
\newcommand{\R}{\mathbb{R}}
\newcommand{\C}{\mathbb{C}}
\newcommand{\F}{\mathbb{F}}
\newcommand{\A}{\mathbb{A}}
\newcommand{\PP}{\mathbb{P}}
\newcommand{\E}{\mathbb{E}}
\newcommand{\eps}{\varepsilon}
\newcommand{\dd}{\,\mathrm{d}}
\begin{document}
\section*{UnsolvedMath ID: 30000370}
\textbf{Source problem code:} catalogue entry
\subsection{Definitions and mathematical statement}
\textbf{Four requests.} Given a family of L\'evy measures \(\mu(x,dh)\) on \(\R^d\), determine when it defines corresponding jump processes; obtain their heat-kernel properties; determine regularity of their harmonic functions; and determine whether the generator extends to that of a Feller semigroup. These are four linked tasks, not a claim that every pointwise family automatically has all the stated properties.\par A L\'evy measure at \(x\) is a positive Borel measure on \(\R^d\setminus\{0\}\) with \(\int(1\wedge|h|^2)\mu(x,dh)<\infty\). To be a jump kernel the family also needs appropriate measurability in \(x\), not merely this pointwise integrability. One standard compensated candidate, with zero drift relative to the displayed truncation, is\[Lf(x)=\int_{\R^d\setminus\{0\}}\bigl(f(x+h)-f(x)-\mathbf1_{\{|h|\le1\}}\nabla f(x)\cdot h\bigr)\,\mu(x,dh),\qquad f\in C_c^\infty(\R^d).\] The source does not supply its generator convention, drift, symmetry, or spatial regularity, so this formula is explicit explanatory notation rather than an asserted model for every intended case. A process associated with \(L\) can be specified through the martingale problem for \(f(X_t)-f(X_0)-\int_0^tLf(X_s)\,ds\), with any lifetime and localization conditions stated. Existence, uniqueness in law, conservativeness, and the strong Markov property are distinct requirements needing hypotheses.\par If a transition semigroup has a density \(P_t(x,dy)=p(t,x,y)\,dy\), its heat-kernel task includes existence and the intended bounds or regularity of \(p\). Such a Lebesgue density is not automatic for arbitrary jump kernels. In an open set \(D\), a probabilistically harmonic function obeys the appropriate exit-time mean-value property \(h(x)=\mathbb E^x h(X_{\tau_U})\) for suitable \(U\Subset D\), whenever the values and expectations are defined; determining its continuity, H\"older regularity, or Harnack properties requires specifying the desired class. A pointwise equation \(Lh=0\) is not imposed on a function outside the operator's domain.\par A Feller semigroup is a strongly continuous positive contraction semigroup on \(C_0(\R^d)\). The extension question asks whether there is such a generator \(G\) with the intended test-function domain included in \(\operatorname{Dom}(G)\) and \(Gf=Lf\) there; it is not automatically a claim that the unclosed operator \(L\) itself is a generator or that its closure is the only extension. The family class, uniform integrability, behavior at infinity, desired estimates, and domain assumptions are not supplied.\par The supplied triage checked 2026-08-21 labels the record partially solved. It reports constructions, heat-kernel estimates, Harnack and harmonic regularity results, and Feller properties for important variable-order classes while saying the full arbitrary-measure program is not settled. Exact theorem hypotheses are absent, and this dated assessment requires independent verification.
\subsection{Short English statement}
Find hypotheses giving a jump process from the measure family, heat-kernel information, harmonic-function regularity, and a Feller-generator extension. Pointwise L\'evy integrability alone does not specify the required spatial measurability, operator convention, or analytic properties.
\subsection{Sources}
\begin{itemize}\item Local UnsolvedMath registry, permanent ID 30000370, code catalogue entry; \emph{Jump Processes from Families of L\'evy Measures}, Oberwolfach Reports Open Problems.\item \emph{Heat Kernels, Stochastic Processes and Functional Inequalities (2006)}, Oberwolfach Reports, DOI 10.4171/owr/2005/54, pp. 3061--3120. \url{https://doi.org/10.4171/owr/2005/54}. Proposed year and year tag: 2005; no proposer named.\item Supplied literature label \emph{original OWR report}: \url{https://doi.org/10.4171/owr/2005/54}.\item Supplied literature label \emph{jump processes and Feller property}: \url{https://www.numdam.org/articles/10.1214/08-AIHP208/}.\item Supplied literature label \emph{harmonic regularity for jump processes}: \url{https://arxiv.org/abs/1109.3678}.\item Literature triage checked 2026-08-21, status \texttt{partially\_solved}; record updated 2026-09-16. The \texttt{corrected\_verified} field records spelling and notation normalization while preserving all four tasks. Verify dated claims independently.\end{itemize}
\end{document}
